\documentclass[11pt,a4paper]{article}
\usepackage[utf8]{inputenc}
\usepackage[margin=1in]{geometry}
\usepackage{amsmath,amssymb}
\usepackage{booktabs}
\usepackage{xcolor}
\usepackage{hyperref}
\usepackage{fancyhdr}
\usepackage{enumitem}

\setlength{\headheight}{14pt}

\hypersetup{
    colorlinks=true,
    linkcolor=blue,
    filecolor=magenta,      
    urlcolor=cyan,
    pdftitle={Project Report: Arduino RFID Access Control System (RC522)},
    pdfpagemode=UseOutlines,
}

\pagestyle{fancy}
\fancyhf{}
\rhead{Arduino RFID Access Control System}
\lhead{Technical Project Report}
\rfoot{Page \thepage}

\title{\textbf{\Huge Engineering Project Report}\\[0.5em] \Large Project 042: Arduino RFID Access Control System (RC522)}
\author{\textbf{Bigyanlabs Engineering Team}\\ Electronics \& Embedded Systems Division}
\date{\today}

\begin{document}

\maketitle
\thispagestyle{empty}

\begin{abstract}
This technical report presents the comprehensive design, circuit schematics, bill of materials, mathematical modeling, assembly methodology, and experimental testing for \textbf{Arduino RFID Access Control System (RC522)}. Classified under the \textbf{Intermediate (College Level)} tier in the \textbf{Security / Embedded Systems} category, this design provides optimized performance, high reliability, and clear practical utility for school and engineering applications.
\end{abstract}

\newpage
\tableofcontents
\newpage

\section{Introduction \& System Overview}
The objective of this project is to implement a robust electronic system for \textbf{Arduino RFID Access Control System (RC522)}. This document details the step-by-step engineering design, component functions, mathematical principles, and experimental results.

\section{Bill of Materials (Components List)}
\begin{table}[h!]
\centering
\caption{Bill of Materials for Arduino RFID Access Control System (RC522)}
\vspace{0.5em}
\begin{tabular}{clcp{7cm}}
\toprule
\textbf{S.No.} & \textbf{Component Name} & \textbf{Qty} & \textbf{Circuit Function} \\
\midrule
1 & Microcontroller / Op-Amp & 1 & Arduino UNO / ESP8266 / ESP32 / LM358 \\
2 & Transducers / Sensors & 2 & Interface sensor modules \\
3 & LCD / OLED Display & 1 & Visual data display (16x2 / I2C) \\
4 & Relays / Driver MOSFETs & 2 & Actuator control interface \\
5 & Passives and Connectors & 1 Set & Resistors, capacitors, breadboard / PCB \\
6 & Power Module & 1 & 5V / 12V regulated DC source \\
\bottomrule
\end{tabular}
\end{table}

\section{System Architecture \& Methodology}
The Arduino RFID Access Control System (RC522) integrates hardware sensors and microcontrollers/op-amps. Sensor signals are processed through custom firmware or analog conditioning stages and displayed locally or streamed over Wi-Fi/Bluetooth.

\section{Circuit Assembly \& Testing Procedure}
\begin{enumerate}
    \item Place core active components and microcontrollers/ICs on the breadboard or PCB.
    \item Wire DC supply rails (+VCC and GND) with appropriate decoupling capacitors.
    \item Connect sensor networks and input signal conditioning circuits.
    \item Interface output switches, MOSFET drivers, or visual/acoustic indicators.
    \item Apply power and measure operating node voltages to verify proper performance.
\end{enumerate}

\section{Mathematical Derivations \& Governing Equations}
\begin{equation}
V_{\text{ADC}} = \frac{\text{ADC\_Value}}{1023} \cdot 5.0\,\text{V}
\end{equation}
\begin{equation}
\text{Distance} = \frac{\text{Time} \times 0.0343\,\text{cm/s}}{2}
\end{equation}

\section{Applications \& Engineering Conclusion}
Engineering laboratory projects, IoT sensor nodes, automated environmental monitoring, and embedded security / embedded systems.

\end{document}
